\section{前言}

结构优化设计旨在在满足强度、刚度、稳定性和制造约束的前提下，合理分配材料、调整结构边界或优化关键尺寸，从而提升结构综合性能。按照设计变量不同，相关研究通常可划分为尺寸优化、形状优化和拓扑优化三类\cite{Wu2021SMO}。传统有限元方法虽然较为成熟，但在复杂曲面边界、薄壁壳体拓扑以及频繁几何更新场景下，几何重建、网格生成和分析模型维护仍然构成主要瓶颈。

等几何分析将 CAD 中常用的 B 样条、NURBS、细分曲面和 T 样条基函数直接用于数值分析，使几何表示与计算模型保持一致，从而减少“设计模型”与“分析模型”之间的转换损失\cite{Wang2021FMEReconstruction,Wen2024CompStructPHT,Liu2018CJSolid}。这种 CAD/CAE 一体化特征使其适合结构优化问题：设计变量可以直接定义在控制点、厚度场或密度场上，高阶连续基函数也有利于曲面边界更新、壳体弯曲分析和灵敏度计算。

近年来，等几何结构优化已从参数化曲线曲面问题扩展到多补丁壳体、自由形变驱动厚度优化以及非结构 T 样条支撑下的形状-拓扑协同优化\cite{Bandara2019ArXiv,Zhao2024EngComput,Wang2025CADA}。本文围绕等几何分析在智能结构优化设计中的方法进展与工程应用展开，重点讨论其在几何建模、拓扑优化、壳结构优化和智能生成闭环中的作用，并归纳现有问题与未来方向。
